\documentclass[10pt]{beamer}

\usetheme{metropolis}

\usepackage{mathpartir}
\usepackage{wrapfig}
\usepackage{scalerel}
\usepackage{minted}
\usepackage{mathtools}
\usepackage{fontawesome} % lock symbol
\usepackage{tikz}
\usepackage{stmaryrd}
\usepackage[bb=boondox]{mathalfa}
\usepackage{pifont}
\usepackage{listings}


\newcommand{\under}[2]{\underset{\raisebox{2mm}{$\scriptstyle #2$}}{#1}}
\newcommand{\inh}[1]{\under{#1}{\uparrow}}
\newcommand{\syn}[1]{\under{#1}{\downarrow}}
\DeclareMathOperator{\llet}{\text{let}}
\DeclareMathOperator{\lin}{\text{in}}
\DeclareMathOperator{\case}{\text{case}}
\DeclareMathOperator{\bbox}{\text{box}}
\DeclareMathOperator{\unbox}{\text{unbox}}
\DeclareMathOperator{\mmod}{\text{mod}}
\DeclareMathOperator{\thunk}{\text{thunk}}
\DeclareMathOperator{\force}{\text{force}}
\DeclareMathOperator{\return}{\text{return}}
\newcommand{\unmod}[3]{\text{let}_{#1} \text{ mod}_{#2}({#3}) \leftarrow}
\newcommand{\at}{\mathbin{\mbox{\small\ensuremath{@}}}}
\DeclarePairedDelimiter\ann{\langle}{\rangle}
\newcommand{\unit}{\mathbb{1}}
\newcommand{\zero}{\mathbb{0}}
\DeclareMathOperator{\inl}{\text{inl}}
\DeclareMathOperator{\inr}{\text{inr}}
\DeclareMathOperator{\suc}{\text{succ}}
\DeclareMathOperator{\absurd}{\text{absurd}}
\DeclareMathOperator{\while}{\text{while}}
\DeclareMathOperator{\fix}{\text{fix}}
\DeclareMathOperator{\primrec}{\text{primrec}}
\newcommand{\nat}{\mathbb{N}}
\newcommand{\tee}{\vdash}
\newcommand{\teev}{\vdash^\mathrm{v}}
\newcommand{\teec}{\vdash^\mathrm{c}}
\newcommand{\hole}{\,?}
\newcommand*{\Ev}{E_\mathrm{v}}
\newcommand*{\Ep}{K}
\DeclareMathOperator{\red}{\rightsquigarrow}
\DeclarePairedDelimiter\floor{\lfloor}{\rfloor}
\DeclarePairedDelimiter\ceil{\lceil}{\rceil}
\DeclareMathOperator{\lock}{\tikz{\node[scale=1, line width=0pt, inner sep=0pt] at (0,0) {\faLock\,};}}
\definecolor{gainsboro}{rgb}{0.86, 0.86, 0.86}
\newcommand{\reducedstrut}{\vrule width 0pt height .9\ht\strutbox depth .9\dp\strutbox\relax}
\newcommand{\gray}[1]{%
  \begingroup
  \setlength{\fboxsep}{0pt}%
  \colorbox{gainsboro}{\reducedstrut$#1$\/}%
  \endgroup
}

\newcommand{\clashes}{\,\text{\Lightning}\,}
\newcommand{\notclashes}{\,{\text{\Lightning}\mkern-8.5mu\backslash}\,}
\DeclareMathOperator{\joinctx}{\mathbin{\fatsemi}}
\DeclareMathOperator{\splitctx}{\curlyveeuparrow}

\newcommand{\one}{\mathbf{1}}
\newcommand{\sem}[1]{\llbracket #1 \rrbracket}

\renewcommand{\qed}{\hfill\resizebox{8px}{8px}{\input{image.pdf_tex}}}

\DeclareMathOperator{\scale}{\mathbin{\triangleleft}}
\DeclareMathOperator{\boxop}{\mathbin{\triangleright}}

\begin{document}

\begin{frame}[fragile]
\frametitle{Background}

\begin{itemize}
\setlength{\itemsep}{1em}
  \item What are modes?
  \begin{lstlisting}[language = caml, basicstyle=\footnotesize\ttfamily, aboveskip=0.5, belowskip=0.5]
let process_parts_wrong (pair @ unique) =
let x @ unique = pair.fst in
let y @ unique = pair.fst in ... 
  \end{lstlisting}
  \item Modes as semirings (with orderings). \\
  many $\le$ once \\
  once + once = many \\
  unique many $\cdot$ aliased once = aliased many \\
  \item Where the semiring goes wrong... \\
$\cdot, y : \tau \at \text{AM} \tee (\bbox_A y, \bbox_A y) : \boxempty^A \tau \times \boxempty^A \tau \at \text{CA}$
$\cdot, y : \tau \at \text{AM} \tee \llet_{\bot}x =  \bbox_A y \lin (x, x) : \boxempty^A \tau \times \boxempty^A \tau \at \text{CA}$ ???
\item Again, what are modes?
\end{itemize}

\end{frame}

\begin{frame}
\vspace{0.7em}
\frametitle{What I did}
Proved useful properties (substitution, beta and eta law type preservation) of various modal calculi to compile a minimal set of laws needed for the modes that we wanted to express. \\

Takeaways:
\begin{itemize}
  \item Splitting $\cdot$ into $\scale$ (scale) and $\boxop$ (box).
  \item Converting equalities to inequalities. \\
  $(a \scale b) \scale c = a \scale \, (b \scale c)$ might not hold, but $(a \scale b) \scale c \le a \scale \, (b \scale c)$ does!
  \item Locks expressed as a division operator in terms of $\scale$.
\end{itemize}

What else you can do:
\begin{itemize}
  \item Show other mode systems follow our structure.
  \item Take advantage of the weaker laws to do other fun things (Kleene star for grades!)
\end{itemize}
\end{frame}

\end{document}